\documentclass{strategynet-report}

\hypersetup{
  pdftitle={Slice + Hypercube II: Circuits and Replay},
  pdfauthor={research@strategynet.ai}
}

\lstdefinestyle{circuit}{
  style=sncode
}

\newcommand{\SLICE}{\textsf{Slice}}
\newcommand{\CIRCUIT}{\textsf{Circuit}}
\newcommand{\R}{\mathbb{R}}

\tikzset{
  circuitbox/.style={
    draw=SNInk,
    fill=SNInk!5,
    rounded corners=2pt,
    align=center,
    minimum height=9mm,
    text width=28mm,
    font=\sffamily\fontsize{8}{9.5}\selectfont
  },
  circuithot/.style={
    circuitbox,
    draw=SNSignalDark,
    fill=SNSignal!7
  },
  circuitmeta/.style={
    circuitbox,
    draw=SNMuted,
    fill=SNMuted!7
  },
  circuitstate/.style={
    circuitbox,
    draw=SNViolet,
    fill=SNViolet!7
  },
  circuiteffect/.style={
    circuitbox,
    draw=SNDanger,
    fill=SNDanger!6
  },
  circuitflow/.style={
    -{Latex[length=2mm]},
    draw=SNInk,
    line width=0.75pt
  },
  circuitrecord/.style={
    -{Latex[length=2mm]},
    draw=SNSignalDark,
    dashed,
    line width=0.75pt
  },
  circuitblocked/.style={
    -{Latex[length=2mm]},
    draw=SNMuted,
    dotted,
    line width=0.75pt
  }
}

\title{\SLICE{} + Hypercube II}
\SNSubtitle{Stateful factor circuits, explicit backpressure, and deterministic replay}
\author{\href{mailto:research@strategynet.ai}{research@strategynet.ai}}
\SNDocumentType{Part II: implementation and evidence}
\SNVersion{Version 0.1}
\SNDate{July 2026}
\SNStatus{Open-source working paper}
\SNShortTitle{Hypercube II: Circuits and Replay}

\begin{document}
\maketitle

\begin{abstract}
The first \SLICE{} + Hypercube paper established a live state plane and a pure,
generation-oriented calculation graph.  It deliberately left stateful
function cells, durable replay, transport, and effects as future boundaries.
This companion paper implements a first stateful callback boundary and a
deterministic recording/replay path.

\texttt{hypercube-circuit} accepts one coherent Hypercube snapshot at a time,
orders it through a single-producer Disruptor ring, and invokes deterministic
stateful processors using generation time rather than processor wall time.
The reference processor turns a calculated cross-section into persistent
threshold state with hysteresis and missing-data invalidation.  Every
generation emits both a complete state cross-section for late subscribers and
sparse transition events for synchronized consumers.

The same path records a versioned manifest, complete Hypercube updates,
optional upstream stream positions, semantic snapshot digests, state, and
transitions.  Replay constructs a fresh engine and fresh processor state,
changes only the execution mode, and compares every generation.  A checked-in
demonstration records 40 generations over 32 entities, recalculates 6,400
factor cells, reconstructs 1,280 trigger-state cells and 78 transitions, and
reports zero divergent generations.  The ETF, pairs, browser, and replay
programs provide an executable application ladder.  The paper does not claim
a production Aeron adapter, raw-feed reconstruction, checkpoint seeking,
latency benchmark, or trading authority; it specifies those boundaries
explicitly.
\end{abstract}

\SNtableofcontents

\part{From Factors to Temporal State}

\section{Continuation from Part I}

Part I separated three responsibilities \cite{partone}: \SLICE{} publishes
typed current vectors, Hypercube evaluates a coherent generation through a
declared graph, and external adapters move or act on the resulting state.
That release could calculate a factor, rank, basket discrepancy, or
cointegrating residual.  It could not remember that a condition had persisted
for three generations, distinguish an entry from an already-active signal, or
reconstruct that temporal state from a recording.

Those omissions become visible in an ordinary factor application.  Let
\(f_{i,g}\) and \(h_{i,g}\) be two calculated factors for entity \(i\) at
generation \(g\).  A linear Hypercube node can calculate a spread
\[
  d_{i,g}=f_{i,g}-h_{i,g}.
\]
A trading, alerting, or monitoring application rarely wants a new action on
every update for which \(d_{i,g}>a\).  It usually wants a lifecycle:
\[
  \text{inactive}\rightarrow\text{qualifying}\rightarrow
  \text{active}\rightarrow\text{exited or invalidated}.
\]
That lifecycle is state across generations, not another cross-sectional
normalization within one generation.

\begin{SNthesis}
The useful product is not merely a stream of factor values.  It is a
reconstructible state transition system whose inputs, logical clock, current
state, sparse changes, overload behavior, and effect boundary are all explicit.
\end{SNthesis}

\section{Application ladder}

The public repository now contains four complementary demonstrations
\cite{repository}:
\begin{center}
\begin{tabularx}{\textwidth}{@{}p{0.18\textwidth}p{0.26\textwidth}X@{}}
\toprule
Demonstration & Implemented operation & Architectural evidence \\
\midrule
browser server &
factor graph over a changing stock cross-section; \SLICE{}, JSON, and SSE
publication &
one calculation can support local memory and remote presentation views \\
ETF monitor &
constituent-axis basket reduction followed by ETF-axis ranking &
domain arithmetic and cross-sectional algebra remain separate, visible stages \\
pairs monitor &
cointegrating residual, model standardization, half-life, and opportunity rank &
model values remain aligned to stable pair entities across generations \\
record/replay &
factor calculation, persistent trigger state, versioned recording, and fresh
reconstruction &
the temporal layer can be verified without invoking a live effect \\
\bottomrule
\end{tabularx}
\end{center}

The demonstrations are deliberately not collapsed into one large application.
ETF and pairs make the upstream calculations readable.  The browser makes
publication visible.  The circuit demonstration isolates temporal state and
replay.  Their shared contract is a coherent Hypercube generation.

\section{Three different kinds of truth}

The architecture distinguishes:
\begin{enumerate}
  \item \textbf{current analytical truth}: the latest coherent values and
        callback state a consumer may read;
  \item \textbf{historical reconstruction truth}: the ordered inputs and
        metadata required to rebuild those values and state;
  \item \textbf{effect authority}: the separate capability to publish,
        persist, alert, or trade.
\end{enumerate}

A ring buffer is not durable history.  A memory-mapped vector is not an event
log.  A replayable decision is not permission to execute it.  Keeping these
distinctions structural avoids two common errors: treating a fast queue as the
system of record, and allowing deterministic calculation code to inherit
irreversible authority.

\part{Circuit Contract}

\section{One coherent generation is one event}

For generation \(g\), the pure engine receives
\[
  U_g=(g,t_g,m_g,R_g,N_g)
\]
and emits a coherent snapshot
\[
  H_g=E(U_g)=(g,t_g,m_g,n_g,V_g,Q_g).
\]
Here \(R_g\) is the complete entity-row cross-section, \(N_g\) is the declared
node graph, \(V_g\) is the ordered calculated value set, and \(Q_g\) is node
status.  The circuit does not subscribe to individual cell writes.  Its event
is the immutable \(H_g\), after the engine has accepted the complete update.

\begin{figure}[htbp]
\centering
\begin{tikzpicture}[node distance=7mm and 7mm]
  \node[circuitmeta] (input) {\textbf{Update \(U_g\)}\\rows, graph, logical time};
  \node[circuithot, right=of input] (engine)
    {\textbf{Hypercube}\\pure generation calculation};
  \node[circuithot, right=of engine] (snapshot)
    {\textbf{Snapshot \(H_g\)}\\coherent cross-section};
  \node[circuitstate, right=of snapshot] (ring)
    {\textbf{Disruptor}\\bounded ordered handoff};
  \node[circuitstate, right=of ring] (processor)
    {\textbf{Processor}\\declared temporal state};
  \draw[circuitflow] (input) -- (engine);
  \draw[circuitflow] (engine) -- (snapshot);
  \draw[circuitflow] (snapshot) -- (ring);
  \draw[circuitflow] (ring) -- (processor);
  \node[circuitmeta, below=10mm of processor] (outputs)
    {\textbf{Output}\\state vector + transitions};
  \draw[circuitflow] (processor) -- (outputs);
  \node[circuitmeta, below=10mm of input] (record)
    {\textbf{Recording}\\input + expected output};
  \draw[circuitrecord] (input) -- (record);
  \draw[circuitrecord] (outputs) -- (record);
\end{tikzpicture}
\caption{Implemented calculation, circuit, and recording boundary.  A callback
never observes a half-built Hypercube generation.}
\label{fig:circuit-path}
\end{figure}

The per-frame context is
\[
  K_g=(q_g,g,t_g,m_g),
\]
where \(q_g\) is the circuit sequence.  Stateful code uses \(g\) and \(t_g\)
as logical coordinates.  Processor scheduling time, operating-system time,
and measured calculation duration are operational observations, not inputs to
the state transition function.

\section{Why a Disruptor is useful here}

The LMAX Disruptor separates event storage, sequencing, producer claims,
consumer progress, and dependency graphs around a preallocated bounded ring
\cite{lmax}.  The Rust \texttt{disruptor} crate implements the same family of
patterns for low-latency inter-thread communication and recommends
nonblocking publication when the application must handle a full ring
explicitly \cite{rustdisruptor}.

The current circuit chooses a single producer because Hypercube generations
already have one ordered owner.  One stateful processor consumes them in
sequence.  The ring offers:
\begin{itemize}
  \item a preallocated bounded handoff rather than unbounded queue growth;
  \item a circuit sequence independent of wall-clock timestamps;
  \item nonblocking submission with a visible \texttt{RingFull} result;
  \item configurable sleep, spin-loop, and busy-spin waiting;
  \item an asynchronous submit/receive API and a lockstep correctness path.
\end{itemize}

It does \emph{not} make an overloaded consumer current.  If producer sequence
\(p\), consumer frontier \(c\), and ring capacity \(B\) satisfy
\[
  p-c\ge B,
\]
the next nonblocking publication is rejected.  That failure is valuable
information: the consumer cannot keep up and latency would otherwise grow
behind the producer.

\section{Backpressure is an application policy}

The circuit never silently overwrites an unprocessed correctness-lane
generation.  On a full ring, an owner must choose one of three policies:
\begin{enumerate}
  \item backpressure the producer and preserve every accepted generation;
  \item retain normalized input durably and catch up from its log position;
  \item coalesce or drop only a separately named best-effort view.
\end{enumerate}

\texttt{CaptureSession} and \texttt{ReplayRunner} use lockstep processing.
The engine does not accept the next update until callback processing and the
recording step complete.  Where an upstream market feed cannot be slowed, the
feed should be archived before generation assembly.  The circuit is then a
deterministic consumer of durable input, not the sole owner of history.

\begin{SNnote}[What the ring solves]
The ring makes ordering, capacity, and consumer progress visible inside one
process.  It does not solve raw-feed durability, cross-host recovery,
multi-stream generation assembly, or external-effect idempotency.
\end{SNnote}

\section{Processor determinism}

A processor is a state transition function
\[
  (S_g,O_g)=P(S_{g-1},K_g,H_g),
\]
where \(S_g\) is private declared state and \(O_g\) is the published
deterministic output.  Replay is meaningful when:
\begin{enumerate}
  \item \(P\) reads only \(S_{g-1}\), \(K_g\), \(H_g\), and immutable
        configuration;
  \item iteration and output ordering are canonical;
  \item missing data and reset behavior are specified;
  \item nondeterministic I/O, random sources, and wall time stay outside \(P\);
  \item irreversible effects consume \(O_g\) through another adapter.
\end{enumerate}

Logical ordering has a longer history than this implementation.  Lamport
showed that event ordering can be represented independently of synchronized
physical clocks \cite{lamport}.  Hypercube does not implement a distributed
Lamport clock, but it applies the relevant separation: generation and circuit
sequence determine processing order, while observation time describes the
input.

\section{Reference threshold state machine}

The first processor watches one calculated node for every entity.  A trigger
specification is
\[
  T=(\mathrm{id},\mathrm{node},a,b,p),
\]
with entry threshold \(a\), exit threshold \(b\), required persistence \(p\),
and strict hysteresis \(a>b\).

For an inactive entity with observed value \(x_{i,g}\), the qualifying count is
\[
  k_{i,g}=
  \begin{cases}
    k_{i,g-1}+1,&x_{i,g}\ge a,\\
    0,&x_{i,g}<a.
  \end{cases}
\]
The state becomes active when \(k_{i,g}\ge p\).  An active state remains active
through the hysteresis band and exits only when
\[
  x_{i,g}\le b.
\]
If the watched node has no value for an active entity, the state is
\emph{invalidated}, not silently exited.  If a qualifying but inactive entity
is missing, its count is removed, so a later value cannot falsely complete a
``consecutive'' run across a gap.

\begin{figure}[htbp]
\centering
\begin{tikzpicture}[node distance=13mm and 18mm]
  \node[circuitmeta] (inactive) {\textbf{Inactive}\\\(k=0\)};
  \node[circuitstate, right=of inactive] (qualifying)
    {\textbf{Qualifying}\\\(0<k<p\)};
  \node[circuithot, right=of qualifying] (active)
    {\textbf{Active}\\entry emitted};
  \draw[circuitflow] (inactive) --
    node[above,font=\scriptsize\sffamily]{\(x\ge a\)} (qualifying);
  \draw[circuitflow] (qualifying) --
    node[above,font=\scriptsize\sffamily]{\(k\ge p\)} (active);
  \draw[circuitflow] (qualifying) to[bend left=33]
    node[below,font=\scriptsize\sffamily]{\(x<a\) or missing} (inactive);
  \draw[circuitflow] (active) to[bend left=32]
    node[above,font=\scriptsize\sffamily]{\(x\le b\): exit} (inactive);
  \draw[circuitblocked] (active) to[loop above]
    node[above,font=\scriptsize\sffamily]{\(x>b\)} (active);
  \node[circuitmeta, below=15mm of active] (invalid)
    {\textbf{Absent}\\invalidation emitted};
  \draw[circuitflow] (active) --
    node[right,font=\scriptsize\sffamily]{missing} (invalid);
\end{tikzpicture}
\caption{Persistent upper-threshold trigger with hysteresis and explicit
missing-data invalidation.}
\label{fig:trigger-state}
\end{figure}

\section{Dense state and sparse transitions}

For every generation, the processor emits:
\[
  Z_g=\{(\mathrm{trigger},\mathrm{node},i,
          \mathrm{active}_{i,g},k_{i,g})\}
\]
and
\[
  \Delta Z_g=\{\mathrm{Entered},\mathrm{Exited},
                \mathrm{Invalidated}\text{ events at }g\}.
\]

\(Z_g\) is the complete ordered trigger-state cross-section.  It maps
naturally to a new \SLICE{} vector and lets a late subscriber bootstrap without
replaying the entire run.  \(\Delta Z_g\) is the sparse edge stream for a
subscriber that already holds \(Z_{g-1}\).  Recording both views also exposes
latent state divergence even when no transition happens in a generation.

This dual output is the bridge from an Excel-like continuously recalculated
factor sheet to an application state machine.  Factor values remain
inspectable, signal lifecycle becomes explicit, and downstream consumers can
choose snapshot or event semantics without recalculating either.

\part{Deterministic Record and Replay}

\section{Choosing the replay boundary}

``Replay the system'' is incomplete until the recorded boundary is named.
\begin{center}
\begin{tabularx}{\textwidth}{@{}p{0.20\textwidth}p{0.27\textwidth}X@{}}
\toprule
Recorded boundary & What replay reconstructs & Tradeoff \\
\midrule
published factor frame &
callbacks, strategies, and downstream state only &
smallest record; cannot verify factor calculation \\
complete Hypercube update &
factor graph, values, callback state, and transitions &
implemented compromise; repeats rows and graph per generation \\
normalized raw events plus generation seals &
ingestion, generation assembly, factors, and callbacks &
strongest apparatus replay; requires ordering, watermark, and join contracts \\
\bottomrule
\end{tabularx}
\end{center}

The current implementation records complete Hypercube updates.  It is more
useful than factor-only replay because changes to rows, node graphs, transforms,
or arithmetic can be detected.  It stops short of claiming that vendor
normalization and multi-feed generation assembly are reproducible.

\section{Versioned recording}

A recording is
\[
  \mathcal D=(M,G_0,G_1,\ldots,G_n).
\]
The manifest
\[
  M=(v,\mathrm{run},\mathrm{build},\mathrm{config},
     \mathrm{layout},\mathrm{triggers},\mathrm{metadata})
\]
identifies the schema and non-generation configuration.  One generation
record is
\[
  G_g=(q_g,U_g,F_g,D_g,Z_g,\Delta Z_g),
\]
where \(F_g\) is an optional map from stable upstream stream identity to the
last position included in the generation and \(D_g\) is the expected semantic
snapshot digest.

The development adapter writes one manifest and then one generation per JSON
line.  It validates:
\begin{itemize}
  \item exactly one leading version-1 manifest;
  \item contiguous circuit sequences and increasing generations;
  \item agreement between update and expected output generation;
  \item valid trigger configuration and nonempty provenance fields;
  \item only finite JSON numeric inputs.
\end{itemize}

JSON parsing enables exact finite \(f64\) round trips.  NaN and infinity are
rejected before the engine advances because JSON would otherwise encode them
as \texttt{null}; a missing primitive is represented by omitting its field.
The format is intentionally inspectable and not presented as the final
low-latency wire encoding.

\section{Semantic exactness}

The semantic snapshot digest includes:
\begin{itemize}
  \item generation, observation time, and entity count;
  \item ordered node and entity identities;
  \item exact IEEE-754 value bits and value observation times;
  \item ordered status identities and value/missing counts.
\end{itemize}

It excludes:
\begin{itemize}
  \item \texttt{ExecutionMode}, because \texttt{Live} becomes
        \texttt{Replay};
  \item measured node compute duration;
  \item wall time, thread identity, transport delay, and logging details.
\end{itemize}

The current digest is a versioned streaming xxHash64 over a canonical byte
representation.  Digest equality is supplemented by exact comparison of the
full ordered state and transition outputs.  The build and configuration
identities remain provenance: a different compiler, target, or implementation
may be replayed, but any arithmetic difference is reported rather than
normalized away.

\begin{SNnote}[Exact does not mean identical runtime]
Replay requires semantic equality of calculated values and state transitions.
It does not require identical processor scheduling, compute duration, or
transport timing.  Including those operational measurements would make a
correct replay fail for irrelevant reasons.
\end{SNnote}

\section{Replay algorithm}

Replay performs the following bounded procedure:
\begin{enumerate}
  \item validate and read the recording manifest;
  \item construct a new \texttt{HypercubeEngine};
  \item construct a new processor from the recorded trigger configuration;
  \item set each recorded update mode to \texttt{Replay};
  \item calculate the snapshot and process it in lockstep;
  \item compare circuit sequence, semantic digest, complete state, and
        transitions;
  \item count every divergent generation and retain the first concrete
        expected-versus-actual difference.
\end{enumerate}

Formally,
\[
  H'_g=E'(U_g[m\leftarrow\mathrm{Replay}]),
\]
\[
  (S'_g,O'_g)=P'(S'_{g-1},K'_g,H'_g).
\]
An exact replay requires
\[
  q'_g=q_g,\qquad D(H'_g)=D_g,\qquad
  Z'_g=Z_g,\qquad\Delta Z'_g=\Delta Z_g
\]
for every recorded generation.

Research on deterministic component replay likewise separates virtual or
logical ordering from incidental runtime scheduling \cite{detreplay}.  The
Hypercube implementation is narrower: it records already coherent
application updates and executes one ordered stateful lane.

\section{Divergence is a result, not an exception}

Malformed files, invalid updates, and processor failures are operational
errors.  A valid recording that produces a different digest, state, or
transition is a semantic result.  The CLI therefore uses separate exits:
\begin{center}
\begin{tabular}{@{}cl@{}}
\toprule
Exit & Meaning \\
\midrule
0 & every recorded generation matched \\
1 & operational or recording-format failure \\
2 & replay completed with semantic divergence \\
\bottomrule
\end{tabular}
\end{center}

This distinction supports two workflows:
\begin{itemize}
  \item \textbf{reconstruction}: expect zero divergence under the same
        implementation contract;
  \item \textbf{counterfactual replay}: intentionally change a node, threshold,
        or model and analyze the resulting divergence.
\end{itemize}

\section{Effects remain outside replay}

Event-sourced systems must prevent replayed events from reissuing external
updates to systems that cannot distinguish replay from live processing
\cite{eventsourcing}.  Hypercube enforces the simpler precursor:
the MVP contains no execution adapter, and processors are documented as
deterministic calculation code without irreversible effects.

Future replay publication should use a namespace such as
\[
  \texttt{replay/\{run\_id\}/\ldots}
\]
and every authorized effect adapter must reject replay-mode inputs.  A live
order, notification, or database mutation should also carry an idempotency key
and an evidence reference.  Calculation replay and effect replay are separate
capabilities.

\part{Transport and Deployment}

\section{Local circuit versus Aeron}

Inside one compute process, Aeron is unnecessary for the engine-to-callback
handoff.  The immutable snapshot is already in memory; the Disruptor provides
ordered inter-thread transfer without serialization or a media driver.

Aeron becomes useful at a process or host boundary:
\begin{itemize}
  \item factor/state publication to independent services;
  \item durable capture of normalized inputs or generation envelopes;
  \item subscriber recovery from a known stream position;
  \item controlled replay onto a dedicated channel.
\end{itemize}

Aeron Archive records and replays streams from durable storage and supports
replay from a specified recording and position \cite{aeron}.  The current
crate does not depend on Aeron.  Its public manifest and generation structs are
transport-neutral envelopes that a future SBE adapter can encode as:
\begin{lstlisting}[style=circuit]
hypercube.recording.manifest.v1
hypercube.recording.generation.v1
\end{lstlisting}

\section{Source positions define the input frontier}

For \(k\) upstream streams, generation \(g\) may record
\[
  F_g=(p_{1,g},p_{2,g},\ldots,p_{k,g}),
\]
where \(p_{j,g}\) is the last input position included from stream \(j\).
The manifest can carry Archive recording identities and start positions;
\texttt{source\_positions} carries the per-generation frontier.

This vector makes a generation's input cut inspectable, but it does not by
itself prove that the cut is semantically coherent.  The application must
still define:
\begin{itemize}
  \item the canonical order when inputs from different streams race;
  \item event-time and ingestion-time watermark rules;
  \item late, duplicate, corrected, and missing event behavior;
  \item the generation seal that says no more input belongs to \(g\).
\end{itemize}

Those rules are the remaining step from complete-update replay to raw-feed
apparatus replay.

\section{Deployment patterns}

\begin{figure}[htbp]
\centering
\resizebox{\textwidth}{!}{%
\begin{tikzpicture}[node distance=7mm and 9mm]
  \node[circuitmeta] (feed) {\textbf{Normalized feed}\\Aeron or local adapter};
  \node[circuitmeta, right=of feed] (assembler)
    {\textbf{Assembler}\\coherent update + frontiers};
  \node[circuithot, right=of assembler] (cube)
    {\textbf{Hypercube}\\factors and scores};
  \node[circuitstate, right=of cube] (circuit)
    {\textbf{Circuit}\\state and transitions};
  \node[circuitmeta, right=of circuit] (views)
    {\textbf{Views}\\Slice, Aeron, API};
  \draw[circuitflow] (feed) -- (assembler);
  \draw[circuitflow] (assembler) -- (cube);
  \draw[circuitflow] (cube) -- (circuit);
  \draw[circuitflow] (circuit) -- (views);
  \node[circuitmeta, below=10mm of assembler] (archive)
    {\textbf{Archive}\\raw or update boundary};
  \draw[circuitrecord] (feed) -- (archive);
  \draw[circuitrecord] (assembler) -- (archive);
  \node[circuiteffect, below=10mm of views] (effect)
    {\textbf{Effect adapter}\\separate authority};
  \draw[circuitblocked] (views) -- (effect);
\end{tikzpicture}%
}
\caption{Target deployment: live calculation (solid), recording (dashed), and
separately authorized effects (dotted).}
\label{fig:deployment}
\end{figure}

The topology supports a \textbf{single process}; a \textbf{split deployment}
that keeps engine-to-circuit transfer local but publishes views over Aeron; or
\textbf{durable reconstruction} that archives normalized input and publishes
replay through an isolated namespace.

\part{Executable Demonstrations}

\section{ETF: two axes remain visible}

The ETF monitor reduces 160 correlated constituent returns into 12 basket fair
values and then evaluates five ETF-aligned Hypercube nodes.  For ETF \(j\),
\[
  r^{\mathrm{NAV}}_{j,g}=\sum_i w_{ji}r_{i,g},\qquad
  V_{j,g}=V_{j,g-1}(1+r^{\mathrm{NAV}}_{j,g}),
\]
while the simulated market premium follows
\[
  u_{j,g}=0.88u_{j,g-1}+0.00040\eta_{j,g}.
\]
The graph publishes fair value, market price, premium basis points, model
z-score, and cross-sectional premium rank.  The terminal therefore displays
60 cells per generation.

\begin{figure}[htbp]
\centering
\includegraphics[width=\textwidth]{etf-demo.png}
\caption{Checked-in ETF recording at generation 28.  Labels preserve the
premium sign independently of color.  Fees, borrow, and creation-unit
frictions remain explicit omissions.}
\label{fig:etf-demo}
\end{figure}

This is an upstream circuit application: a callback may watch
\texttt{premium\_z}, require persistence, and publish entry/exit state without
changing the basket calculation.  A more complete application would combine
factor state with inventory, executable spread, creation size, and policy
before producing a bounded intent.

\section{Pairs: model residual becomes signal state}

The pairs monitor publishes 120 cells over 24 stable pair entities.  Its
residual and standardization are
\[
  y_{j,g}=\log A_{j,g}-\alpha_j-\beta_j\log B_{j,g},
\qquad
  z_{j,g}=\frac{y_{j,g}}{\sigma_j/\sqrt{1-\phi_j^2}}.
\]
Half-life and cross-sectional opportunity rank remain separate nodes.

\begin{figure}[htbp]
\centering
\includegraphics[width=\textwidth]{pairs-demo.png}
\caption{Checked-in pairs recording at generation 28.  The display ranks
standardized dislocations while retaining hedge ratio and half-life context.}
\label{fig:pairs-demo}
\end{figure}

A threshold processor turns a repeated \(z\)-score condition into lifecycle
state.  Hysteresis avoids an exit/re-entry loop near the entry boundary;
invalidation distinguishes a missing leg or suppressed residual from a
genuine convergence exit.  Estimation windows, structural-break detection,
borrow, cost, and portfolio coupling remain outside the demo.

\section{Browser: one generation, several views}

The synthetic server computes five nodes for a changing stock cross-section,
publishes them to entity-aligned \SLICE{} files, exposes a JSON snapshot, and
streams generations through SSE to a dependency-free dashboard.  It
demonstrates that current-state projection and subscriber presentation need
not rerun factor arithmetic.

The circuit extends that fan-out with two additional views:
\begin{itemize}
  \item a dense current-state projection suitable for \SLICE{} or a bootstrap
        response;
  \item a sparse transition stream suitable for already synchronized
        subscribers.
\end{itemize}
Neither browser SSE nor \SLICE{} is the replay log.  They are projections of a
generation whose reproducible input lives at the recording boundary.

\section{Record/replay experiment}

The executable experiment is:
\begin{lstlisting}[style=circuit]
cargo run -p hypercube-circuit --bin hypercube-replay -- \
  record-demo /tmp/hypercube-factor.jsonl 40 32

cargo run -p hypercube-circuit --bin hypercube-replay -- \
  verify /tmp/hypercube-factor.jsonl
\end{lstlisting}

The synthetic graph has five nodes over 32 entities for 40 generations:
\[
  5\cdot32\cdot40=6{,}400
\]
factor cells.  One threshold specification yields
\[
  32\cdot40=1{,}280
\]
complete trigger-state cells.

\begin{center}
\begin{tabularx}{0.82\textwidth}{@{}X r r@{}}
\toprule
Compared object & Recorded & Replayed \\
\midrule
generations & 40 & 40 \\
factor cells represented by semantic digests & 6,400 & 6,400 \\
trigger-state cells & 1,280 & 1,280 \\
transitions & 78 & 78 \\
divergent generations & \multicolumn{2}{r}{0} \\
\bottomrule
\end{tabularx}
\end{center}

The file contains 41 lines: one manifest and 40 generation records.  This is a
functional exactness result, not a throughput or storage benchmark.

\section{Failure contracts exercised by tests}

The circuit tests do more than assert the happy path:
\begin{itemize}
  \item a blocked consumer fills a two-slot ring and the third generation
        returns \texttt{RingFull} without overwriting pending work;
  \item a missing inactive candidate resets persistence, while a missing active
        entity emits \texttt{Invalidated};
  \item sequence gaps and non-increasing generations are rejected;
  \item non-finite JSON input is rejected before a generation is written;
  \item live and replay JSON values reproduce exact floating-point bits;
  \item a changed recorded primitive produces a snapshot-digest divergence;
  \item unchanged input reconstructs snapshot, state, and transitions.
\end{itemize}

These tests target the places where a superficially working event system can
quietly drift: overload, gaps, missing observations, serialization, hidden
state, and nondeterministic comparison.

\section{Current conformance result}

The workspace currently runs 29 tests:
\begin{center}
\begin{tabularx}{0.86\textwidth}{@{}X r l@{}}
\toprule
Suite & Tests & Result \\
\midrule
engine, publisher, and synthetic injectors & 6 & passed \\
ETF and pairs example invariants & 2 & passed \\
\SLICE{} layout, mmap, records, and vector algebra & 7 & passed \\
circuit, triggers, recording, digest, and replay & 14 & passed \\
\midrule
\textbf{Total} & \textbf{29} & \textbf{passed} \\
\bottomrule
\end{tabularx}
\end{center}

The workspace also passes formatting, clippy with warnings denied, and public
API documentation with warnings denied.  These are conformance and packaging
checks; the repository still makes no latency or throughput claim.

\part{Application and Research Sequence}

\section{The business application}

The immediate application is a factor-to-signal state service:
\begin{enumerate}
  \item Hypercube calculates raw fields, factor spreads, ranks, residuals, and
        scores over the full current universe;
  \item the circuit applies persistence, hysteresis, invalidation, and later
        stateful model callbacks;
  \item dense state is published for current consumers;
  \item sparse transitions feed alerts, strategy coordinators, portfolio
        construction, or a separately governed trading adapter;
  \item the recording reconstructs why a state existed and when it changed.
\end{enumerate}

The commercial value is less about inventing another formula and more about
making calculated state reusable.  A factor team owns one factor definition;
multiple applications can consume the same current vector and transition
history.  A strategy can be tested against the exact inputs that drove the
live system.  An operator can distinguish ``no longer qualifies'' from
``source disappeared.''  A new subscriber can bootstrap from dense state
instead of inferring it from an incomplete recent event window.

\section{Counterfactual replay}

Once reconstruction is exact, replay can branch:
\begin{itemize}
  \item change entry, exit, or persistence parameters;
  \item replace one factor node or model version;
  \item compare old and new state transitions by generation;
  \item inject delayed or missing data according to a declared scenario;
  \item measure proposed decision intents without enabling live effects.
\end{itemize}

This resembles the ``parallel model'' enabled by event sourcing: alternative
state can be produced from an alternative event or processing history
\cite{eventsourcing}.  The current CLI verifies only the recorded
configuration.  A counterfactual runner and comparative report are a natural
next layer.

\section{Implemented limits}

The current release establishes a vertical circuit, not a complete event
platform:
\begin{itemize}
  \item JSON Lines is the only recording adapter; Aeron Archive and SBE mapping
        are specified but not implemented;
  \item complete updates are replayed, not raw vendor events or generation
        assembly;
  \item one threshold processor is provided; general callback composition,
        window state, and model artifacts need additional schemas;
  \item replay begins at generation zero; state checkpoints and seek are not
        implemented;
  \item xxHash64 detects ordinary divergence but is not a cryptographic
        authenticity proof;
  \item state vectors are public structs but no dedicated trigger-to-\SLICE{}
        publisher is included;
  \item no benchmark justifies a wait strategy, ring size, or deployment
        topology for production;
  \item no external-effect or trading authority exists in the circuit crate.
\end{itemize}

These limits are intentional.  They keep a replay verification command from
being confused with a production recovery plan or a calculated signal from
being confused with an executable order.

\section{Next implementation sequence}

A bounded continuation is:
\begin{enumerate}
  \item add an SBE codec for manifest and generation envelopes;
  \item record and replay those envelopes through Aeron Archive by identity and
        position;
  \item define canonical multi-stream generation seals and late-event policy;
  \item add periodic processor checkpoints with digest-verified seek;
  \item publish trigger state into \SLICE{} and transitions onto a named
        subscriber channel;
  \item add counterfactual configuration and first-divergence comparison;
  \item benchmark declared topologies with percentile latency, backlog, CPU,
        allocation, and recovery-time measurements;
  \item specify an idempotent simulated effect adapter before any live
        execution integration.
\end{enumerate}

\section{Conclusion}

Part I made current analytical values cheap to share and deterministic to
recalculate within one generation.  Part II adds the missing temporal
boundary.  A coherent snapshot becomes one ordered circuit event; a declared
processor turns it into reconstructible state; a bounded ring exposes overload
instead of hiding it; and a versioned recording proves whether a fresh engine
and processor arrive at the same result.

The architecture now supports a useful end-to-end statement:
\[
  \text{input generation}
  \longrightarrow\text{factor state}
  \longrightarrow\text{signal state}
  \longrightarrow\text{recorded evidence}.
\]
ETF and pairs show readable financial calculations, the browser shows
multi-view publication, and the replay demo verifies the stateful layer.
Aeron, raw-feed reconstruction, checkpoints, benchmarks, and effects remain
separate work, but their contracts are now narrower and testable.

The central design choice is restraint.  The Disruptor moves coherent state;
it is not truth.  The recording reconstructs calculation; it is not authority.
The callback emits a decision state; it is not an effect.  Those distinctions
allow the system to become faster and more composable without making its
behavior less explainable.

\begin{thebibliography}{9}

\bibitem{partone}
StrategyNet,
\emph{\SLICE{} + Hypercube: Current analytical state without rebuilding it
for every reader},
Open-source working paper, version 0.2, July 2026,
\url{https://github.com/ingvor-restrat/hypercube/blob/main/docs/latex/hypercube_foundations.pdf}.

\bibitem{lmax}
M. Thompson, D. Farley, M. Barker, P. Gee, and A. Stewart,
\emph{LMAX Disruptor: High performance alternative to bounded queues for
exchanging data between concurrent threads},
LMAX Exchange, May 2011,
\url{https://lmax-exchange.github.io/disruptor/disruptor.html}.

\bibitem{rustdisruptor}
\emph{disruptor 4.3.0: Low latency inter-thread communication via a ring
buffer},
Rust crate documentation,
\url{https://docs.rs/disruptor/4.3.0/disruptor/}.

\bibitem{lamport}
L. Lamport,
\emph{Time, Clocks, and the Ordering of Events in a Distributed System},
Communications of the ACM, 21(7), pp.~558--565, 1978,
\url{https://www.microsoft.com/en-us/research/publication/time-clocks-ordering-events-distributed-system/}.

\bibitem{eventsourcing}
M. Fowler,
\emph{Event Sourcing},
Further Enterprise Application Architecture, 2005,
\url{https://martinfowler.com/eaaDev/EventSourcing.html}.

\bibitem{detreplay}
R. Strom, C. Dorai, T. H. Feng, and Z. Wei,
\emph{Deterministic Replay for Transparent Recovery in Component-Oriented
Middleware},
IEEE ICDCS, 2009,
\url{https://doi.org/10.1109/ICDCS.2009.79}.

\bibitem{aeron}
Aeron Project,
\emph{Aeron Archive Overview},
\url{https://aeron.io/docs/aeron-archive/overview/}.

\bibitem{repository}
StrategyNet,
\emph{Hypercube public repository},
\url{https://github.com/ingvor-restrat/hypercube}.

\end{thebibliography}

\end{document}
